% =============================================================================
%  paper/supplementary/S2_algorithm_pseudocode.tex
%  Supplementary section S2 -- Pseudocode for six shipped R18
%  algorithms: M22 SE(3) SB solver (W2.1); M23 Transport Relaxation
%  (W2.2); M24 Multi-PH filtration (W2.3); M25 Hilbert primal-dual
%  (unifies W2.1-W2.3); M26 Score-Repellent MC (Wf-9); M28
%  Variance-Controlled PUCT MCTS (Wf-11 / TODO-30).  Each block
%  lists shipped module path, input/output shape contract, per-step
%  cost, and integration point.  Status: SHIPPED (Wf-26 Phase 1,
%  2026-09-26).
%
%  Standalone preamble: defines the honest-framing macros locally so
%  this file compiles as a separate PDF (Digital Discovery SI portal).
% =============================================================================

\section*{Supplementary S2: Algorithm pseudocode}
\addcontentsline{toc}{section}{Supplementary S2: Algorithm pseudocode}

\providecommand{\MEASURED}{\textbf{MEASURED}}
\providecommand{\PROJECTED}{\textbf{PROJECTED}}
\providecommand{\DESIGN}{\textbf{DESIGN}}

\paragraph{Scope.}
Six algorithm blocks in order M22 -- M28.  Each is tagged
\DESIGN{} (algorithm specification shipped as a Python module);
empirical runtimes are \DESIGN{}-with-\MEASURED-unit-tests; the
live training-loop measurements (where available) are
\PROJECTED{}.

% ===========================================================================
\subsection*{M22 -- SE(3) Schroedinger-Bridge solver}
\label{supp:s2-m22}

\paragraph{Source.}
\texttt{molmetal/math\_skills/se3\_schrodinger\_bridge.py},
classes \texttt{SE3Geodesic} / \texttt{LogSinkhornLoss} /
\texttt{SchrodingerBridgeSE3}; adapter at
\texttt{molmetal/adapters/se3\_sb\_integration.py}.

\paragraph{Input.}
Two empirical distributions $\mu_{0}, \mu_{1}$ on $\mathrm{SE}(3)$
as point-clouds $\{g_{0}^{(i)}\}_{i=1}^{N} \subset \mathrm{SE}(3)^{N}$
and $\{g_{1}^{(j)}\}_{j=1}^{M}$; entropy $\varepsilon > 0$;
Sinkhorn iters $T$; geodesic samples
$\{t_{k}\}_{k=0}^{K} \subset [0, 1]$.

\paragraph{Output.}
SB coupling $\Pi^{\star} \in \mathbb{R}^{N \times M}_{+}$; per-pair
geodesic interpolants $\{g^{(i,j)}_{t_{k}}\}_{i, j, k}$.

\begin{algorithmic}[1]
  \STATE \textbf{Input:} source $\{g_{0}^{(i)}\}_{i=1}^{N}$, target
         $\{g_{1}^{(j)}\}_{j=1}^{M}$, $\varepsilon > 0$, $T$, $K$.
  \STATE \textbf{Output:} $\Pi^{\star} \in \mathbb{R}^{N \times M}_{+}$,
         $\{g^{(i,j)}_{t_{k}}\}$.
  \STATE // -- Build the cost matrix on $\mathrm{SE}(3)$ --
  \FOR{$i = 1, \ldots, N$}
    \FOR{$j = 1, \ldots, M$}
      \STATE $\xi \gets \log\!\bigl( (g_{0}^{(i)})^{-1} g_{1}^{(j)} \bigr) \in \mathfrak{se}(3)$
      \STATE $C_{ij} \gets \| \xi \|_{\mathfrak{se}(3)}^{2}$
             \COMMENT{bi-invariant Killing-form norm}
    \ENDFOR
  \ENDFOR
  \STATE // -- Log-domain Sinkhorn fixed-point --
  \STATE $u \gets \mathbf{0} \in \mathbb{R}^{N}$,
         $v \gets \mathbf{0} \in \mathbb{R}^{M}$,
         $K \gets \exp(-C / \varepsilon)$ elementwise
  \FOR{$\tau = 1, \ldots, T$}
    \STATE $u \gets \log a - \mathrm{logsumexp}_{j}
                  \!\bigl( \log K + v^{\top} \bigr)$
    \STATE $v \gets \log b - \mathrm{logsumexp}_{i}
                  \!\bigl( \log K + u \bigr)$
  \ENDFOR
  \STATE $\Pi^{\star} \gets \mathrm{diag}(e^{u}) \, K \, \mathrm{diag}(e^{v})$
  \STATE // -- Walk the SE(3) geodesics --
  \FOR{$i, j$}
    \STATE $\xi_{ij} \gets \log\!\bigl( (g_{0}^{(i)})^{-1} g_{1}^{(j)} \bigr)$
    \FOR{$k = 0, \ldots, K$}
      \STATE $g^{(i,j)}_{t_{k}} \gets \exp\!\bigl( g_{0}^{(i)}, t_{k}\, \xi_{ij} \bigr)$
    \ENDFOR
  \ENDFOR
  \STATE \textbf{return} $\Pi^{\star}$, $\{g^{(i,j)}_{t_{k}}\}$
\end{algorithmic}

\paragraph{Complexity.}
Cost-matrix build $O(NM)$; log-domain Sinkhorn $O(T \cdot NM)$
with $T \approx \|C\|_{\infty} / \varepsilon$ (Altschuler 2017);
geodesic walk $O(NMK)$.  Per Theorem 1 (Mahmood 2026),
$\Pi^{\star}_{\varepsilon} \to \Pi^{\star}$ weakly as
$\varepsilon \to 0$.

\paragraph{Integration point.}
\texttt{integrate\_with\_cfm} (W2.1 adapter) consumes the
geodesics as a drop-in replacement for the Euclidean CFM
interpolant $X_{t} = (1 - t) X_{0} + t X_{1}$; shape contract is
$(B, K, 4, 4)$ for $(B \cdot N \cdot M)$ rigid-motion matrices at
$K$ time steps, matching
\texttt{flow\_matching.solver.ode\_solver.ODESolver}.

% ===========================================================================
\subsection*{M23 -- Transport Relaxation for reward shaping}
\label{supp:s2-m23}

\paragraph{Source.}
\texttt{molmetal/math\_skills/transport\_relaxation.py}, classes
\texttt{RelaxedOT} / \texttt{BoltzmannEndpoint} /
\texttt{TransportRelaxationTerm}; adapter at
\texttt{molmetal/adapters/transport\_relaxation\_integration.py}.

\paragraph{Input.}
Marginals $a \in \Delta_{N}$, $b \in \Delta_{M}$; cost
$C \in \mathbb{R}^{N \times M}_{+}$; $\lambda_{\mathrm{relax}} \ge 0$;
reference Boltzmann $q_{T}$; Sinkhorn iters $T$.

\paragraph{Output.}
Relaxed coupling $\tilde{\Pi}^{\star}_{\varepsilon}$; transport-
relaxation loss term
$\lambda_{\mathrm{relax}}\, \mathrm{KL}(p_{T} \| q_{T})$.

\begin{algorithmic}[1]
  \STATE \textbf{Input:} $a$, $b$, $C$, $\lambda_{\mathrm{relax}}$,
         $\varepsilon$, $q_{T}$, $T$.
  \STATE \textbf{Output:} $\tilde{\Pi}^{\star}_{\varepsilon}$,
         $\mathrm{loss} = \lambda_{\mathrm{relax}}\, \mathrm{KL}(p_{T} \| q_{T})$.
  \STATE // -- Build soft marginals --
  \STATE $\alpha \gets \lambda_{\mathrm{relax}} / (1 + \lambda_{\mathrm{relax}})$
  \STATE $\tilde{a} \gets (1 - \alpha)\, a + \alpha\, \mathbb{1}_{N} / N$
  \STATE $\tilde{b} \gets (1 - \alpha)\, b + \alpha\, \mathbb{1}_{M} / M$
  \STATE // -- Log-domain Sinkhorn on soft marginals --
  \STATE $u \gets \log \tilde{a}$, $v \gets \log \tilde{b}$, $K \gets \exp(-C / \varepsilon)$
  \FOR{$\tau = 1, \ldots, T$}
    \STATE $u \gets \log \tilde{a} - \mathrm{logsumexp}_{j}(\log K + v^{\top})$
    \STATE $v \gets \log \tilde{b} - \mathrm{logsumexp}_{i}(\log K + u)$
  \ENDFOR
  \STATE $\tilde{\Pi}^{\star}_{\varepsilon} \gets \mathrm{diag}(e^{u}) K \mathrm{diag}(e^{v})$
  \STATE // -- Transport-relaxation KL term --
  \STATE $p_{T} \gets \tilde{\Pi}^{\star}_{\varepsilon} \mathbf{1}_{M}$
  \STATE $\mathrm{loss} \gets \lambda_{\mathrm{relax}}\, \mathrm{KL}(p_{T} \| q_{T})$
  \STATE \textbf{return} $\tilde{\Pi}^{\star}_{\varepsilon}$, $\mathrm{loss}$
\end{algorithmic}

\paragraph{Complexity.}
Cost-matrix build $O(NM)$; Sinkhorn $O(T \cdot NM)$; total wall-
clock $O(T \cdot NM)$.  Per Theorem 2 (Jiang-Xu-Zhang 2026),
$|S_{\varepsilon}(\tilde{a}, \tilde{b}; C) - S_{\varepsilon}(a, b;
C)| \le C_{\mathrm{Lip}}\, \lambda_{\mathrm{relax}}\, \varepsilon
\, \|C\|_{\infty}$.

\paragraph{Integration point.}
\texttt{augment\_learned\_shaping} (line $\sim$140 of
\texttt{molmetal/adapters/transport\_relaxation\_integration.py})
appends $\mathrm{loss}$ to the
\texttt{molmetal/molmetal\_lam/reward/learned\_shaping.py} shaping
term.  The wrapper does not modify the live
\texttt{learned\_shaping.py} (R17 close-out is locked).

% ===========================================================================
\subsection*{M24 -- Multi-PH filtration for metal coordination}
\label{supp:s2-m24}

\paragraph{Source.}
\texttt{molmetal/math\_skills/multi\_parameter\_ph.py}, functions
\texttt{two\_param\_filtration} / \texttt{TwoParameterDiagram} /
\texttt{sliced\_wasserstein\_kernel}; adapter at
\texttt{molmetal/adapters/multi\_ph\_integration.py}; guardrail
at \texttt{molmetal/math\_skills/ph\_hardness\_guardrail.py}.

\paragraph{Input.}
Pocket point cloud $\{(x_{i}, z_{i})\}_{i=1}^{N}$; grid
$n_{\mathrm{bins}}$; SW projections $K$.

\paragraph{Output.}
2-PH feature map $F \in \mathbb{R}^{n_{\mathrm{bins}} \times
n_{\mathrm{bins}}}$; SW-kernel $k \in \mathbb{R}^{B \times B}_{+}$.

\begin{algorithmic}[1]
  \STATE \textbf{Input:} pocket $\{(x_{i}, z_{i})\}$, $n_{\mathrm{bins}}$, $K$.
  \STATE \textbf{Output:} $F \in \mathbb{R}^{n_{\mathrm{bins}}^{2}}$, $k \in \mathbb{R}^{B \times B}$.
  \STATE // -- Build the (r, d) edge set --
  \STATE $\mathcal{E} \gets \{ (i, j) : \| x_{i} - x_{j} \| \le r_{\max}
                              \;\wedge\; |z_{i} - z_{j}| \le d_{\max} \}$
  \STATE // -- Fourier-accelerated edge-count (Theorem 1 Kim-Jung 2026) --
  \STATE $E(r, d) \gets \mathrm{FFT}^{-1}\!\bigl(
                \mathrm{FFT}(I_{r}) \cdot \mathrm{FFT}(I_{d}) \bigr)$
         \COMMENT{$O(N \log N)$}
  \STATE // -- Discretise onto the grid --
  \FOR{$\rho, \delta = 1, \ldots, n_{\mathrm{bins}}$}
    \STATE $r \gets r_{\min} + (\rho - 1)\, \Delta r$, $d \gets d_{\min} + (\delta - 1)\, \Delta d$
    \STATE $\mathcal{K}_{r, d} \gets \mathrm{VR}_{\mathcal{E}}(r, d)$
    \STATE $F_{\rho, \delta} \gets |H_{0}(\mathcal{K}_{r, d})| +
                                    |H_{1}(\mathcal{K}_{r, d})|$
  \ENDFOR
  \STATE // -- Sliced-Wasserstein kernel --
  \FOR{$b, b' = 1, \ldots, B$}
    \STATE $\mathrm{sw}_{b, b'} \gets 0$
    \FOR{$k = 1, \ldots, K$}
      \STATE $\theta_{k} \gets 2\pi\, k / K$
      \STATE $\mathrm{sw}_{b, b'} \gets \mathrm{sw}_{b, b'} +
                  W_{1}(\mathrm{proj}_{k}(F_{b}),
                       \mathrm{proj}_{k}(F_{b'}))^{2}$
    \ENDFOR
    \STATE $k_{b, b'} \gets \exp(-\tau^{2}\, \mathrm{sw}_{b, b'} / K)$
  \ENDFOR
  \STATE \textbf{return} $F$, $k$
\end{algorithmic}

\paragraph{Complexity.}
Fourier-accelerated edge-count $O(N \log N)$; grid discretisation
$O(n_{\mathrm{bins}}^{2}\, N)$ dominated by the Vietoris--Rips
filtration per cell; SW-kernel $O(K B^{2}\, n_{\mathrm{bins}}^{2})$.
The Kolbe--Mayr 2026 guardrail caps the per-cell Vietoris--Rips
compute at $200$ off-diagonal pairs and switches to the
$O(n \log n)$ \texttt{approximate\_bottleneck} matcher beyond that.

\paragraph{Integration point.}
The W2.3 adapter \texttt{MultiPHPocketFeatures} (line $\sim$180 of
\texttt{molmetal/adapters/multi\_ph\_integration.py}) consumes $F$
as a $(B, n_{\mathrm{bins}}, n_{\mathrm{bins}})$ tensor and feeds
it to \texttt{MultiPHMetalCoordHead}; \texttt{augment\_pocket\_embedding}
(line $\sim$880) is the drop-in replacement for the
single-parameter Betti-curve summary.

% ===========================================================================
\subsection*{M25 -- Hilbert-space primal-dual adapter}
\label{supp:s2-m25}

\paragraph{Source.}
Unifying operator-theoretic frame of \S\ref{supp:s1-hilbert},
instantiated as the W2.1 / W2.2 / W2.3 adapters jointly.  The
pseudocode below is the \emph{generic} OTPD instance; the three
W2 adapters differ only in the choice of $K$, $F$, $G$.

\paragraph{Input.}
$K : \mathcal{H}_{1} \to \mathcal{H}_{2}$; $G : \mathcal{H}_{1} \to
\mathbb{R} \cup \{+\infty\}$; $F^{\ast} : \mathcal{H}_{2} \to
\mathbb{R} \cup \{+\infty\}$; step sizes $\sigma, \tau > 0$
satisfying $\sigma \tau \| K \|^{2} < 1$; max iters $T_{\max}$.

\paragraph{Output.}
Saddle point $(x^{\star}, y^{\star})$ of
$\min_{x} G(x) + \langle Kx, y \rangle - F^{\ast}(y)$.

\begin{algorithmic}[1]
  \STATE \textbf{Input:} $K$, $G$, $F^{\ast}$, $\sigma$, $\tau$, $T_{\max}$.
  \STATE \textbf{Output:} $(x^{\star}, y^{\star})$.
  \STATE $x_{0} \gets \mathbf{0} \in \mathcal{H}_{1}$, $y_{0} \gets \mathbf{0} \in \mathcal{H}_{2}$
  \FOR{$k = 1, \ldots, T_{\max}$}
    \STATE $y_{k} \gets \mathrm{prox}_{\sigma F^{\ast}}(y_{k-1} + \sigma K x_{k-1})$
    \STATE $x_{k} \gets \mathrm{prox}_{\tau G}(x_{k-1} - \tau K^{\ast} y_{k})$
    \STATE $\bar{x}_{k} \gets \tfrac{1}{k} \sum_{i=1}^{k} x_{i}$,
           $\bar{y}_{k} \gets \tfrac{1}{k} \sum_{i=1}^{k} y_{i}$
  \ENDFOR
  \STATE \textbf{return} $\bar{x}_{T_{\max}}$, $\bar{y}_{T_{\max}}$
\end{algorithmic}

\paragraph{Complexity.}
Per iteration: one $\mathrm{prox}$ evaluation per functional plus
one $K$, $K^{\ast}$ application; ergodic average $O(1)$ per step.
Convergence rate $O(1/k)$ on the primal-dual gap under
$\sigma \tau \|K\|^{2} < 1$ (Chambolle--Pock 2011; Sun--Zheng 2026
Theorem 3.1).

\paragraph{Instantiation table.}
\begin{itemize}
  \item \textbf{W2.1 (M22):}  $K$ = SE(3) marginal-projection;
        $G = \langle C, P \rangle + \varepsilon
        \mathrm{KL}(P \| \mu_{0} \otimes \mu_{1})$;
        $F^{\ast} = \delta_{\{ P \mathbf{1} = a \}}$.
  \item \textbf{W2.2 (M23):}  $K$ = relaxed-marginal projection;
        $G$ as above; $F^{\ast} = \lambda_{\mathrm{relax}}\,
        \mathrm{KL}(\pi_{T} \| q_{T})$.
  \item \textbf{W2.3 (M24):}  $K$ = 2-PH feature-map extraction;
        $G$ = SW-kernel cost; $F^{\ast}$ = SW-kernel positivity
        cone indicator.
\end{itemize}

% ===========================================================================
\subsection*{M26 -- Score-Repellent Monte Carlo}
\label{supp:s2-m26}

\paragraph{Source.}
\texttt{molmetal/molmetal\_lam/sbdd\_env/score\_repellent\_mc.py}
(Wf-9 ship).

\paragraph{Input.}
Score function $s_{\theta} : \mathbb{R}^{d} \to \mathbb{R}^{d}$;
step size $\eta > 0$; repulsive scale $L > 0$; mixing weight
$\alpha \in [0, 1]$; chain length $T$.

\paragraph{Output.}
Markov chain $(x_{0}, x_{1}, \ldots, x_{T})$ targeting
$\rho \propto \exp(-U)$ with $U$ a learned potential.

\begin{algorithmic}[1]
  \STATE \textbf{Input:} $s_{\theta}$, $\eta$, $L$, $\alpha$, $T$.
  \STATE \textbf{Output:} chain $(x_{0}, \ldots, x_{T})$.
  \STATE $x_{0} \sim \rho_{0}$
  \FOR{$t = 1, \ldots, T$}
    \STATE // Langevin step
    \STATE $\xi_{t} \sim \mathcal{N}(0, 2 \eta I)$
    \STATE $x_{t}^{\mathrm{L}} \gets x_{t-1} + \eta\, s_{\theta}(x_{t-1}) + \xi_{t}$
    \STATE // Repulsive correction (Burch--Bardsley 2022)
    \STATE $g \gets s_{\theta}(x_{t}^{\mathrm{L}}) - s_{\theta}(x_{t-1})$
    \STATE $x_{t}^{\mathrm{R}} \gets x_{t}^{\mathrm{L}} + \alpha\, L\, (g / \|g\|)$
    \STATE // MH acceptance
    \STATE $\log \rho_{t} \gets \log \rho_{\theta}(x_{t}^{\mathrm{R}}) -
                            \log \rho_{\theta}(x_{t-1})$
    \STATE $u \sim \mathrm{Uniform}(0, 1)$
    \IF{$\log u < \log \rho_{t}$}
      \STATE $x_{t} \gets x_{t}^{\mathrm{R}}$
    \ELSE
      \STATE $x_{t} \gets x_{t-1}$
    \ENDIF
  \ENDFOR
  \STATE \textbf{return} $(x_{0}, \ldots, x_{T})$
\end{algorithmic}

\paragraph{Complexity.}
Each step is one EGNN forward pass plus one MH correction; per
Burch--Bardsley 2022 \S 4.4 and R18 \S 6.2,
$\mathrm{Var}(\mathrm{score\_repellent}) \approx (1 - \alpha)\,
\mathrm{Var}(\mathrm{Dirichlet})$ for $\alpha \approx 0.2$ in the
canonical $\eta = 0.1$, $L = 1.0$ regime.

\paragraph{Integration point.}
Score-Repellent MC is the inner-loop sampler of the R18 MCTS
rollout: each leaf is sampled by a $T$-step walk, scored by the
seven-channel \texttt{RewardAggregator}, and fed to the SynFlowNet
advantage estimator.  The A/B harness
\texttt{compare\_variance\_vs\_dirichlet\_mcts} (line $\sim$240 of
\texttt{variance\_mcts\_integration.py}) reports the variance
ratio.

% ===========================================================================
\subsection*{M28 -- Variance-Controlled PUCT MCTS}
\label{supp:s2-m28}

\paragraph{Source.}
\texttt{molmetal/molmetal\_lam/search\_alg/
variance\_puct\_integration.py} (Wf-11 ship) +
\texttt{molmetal/molmetal\_lam/search\_alg/
variance\_mcts\_integration.py} (TODO-30 P5.2 ship).

\paragraph{Input.}
PUCT node $s$ with prior $P(\cdot \mid s)$; per-arm statistics
$\{(Q(s, a), N(s, a), \sigma(s, a))\}$; $c_{\mathrm{puct}}$;
$c_{\mathrm{var}}$; $n_{\mathrm{min}}^{\mathrm{visits}}$.

\paragraph{Output.}
$a^{\star} \in \arg\max_{a} \mathrm{PUCT}_{\mathrm{var}}(s, a)$.

\begin{algorithmic}[1]
  \STATE \textbf{Input:} $s$, $P(\cdot \mid s)$, per-arm stats,
         $c_{\mathrm{puct}}$, $c_{\mathrm{var}}$, $n_{\mathrm{min}}^{\mathrm{visits}}$.
  \STATE \textbf{Output:} $a^{\star} \in \mathcal{A}(s)$.
  \STATE $N(s) \gets \sum_{a} N(s, a)$
  \FOR{$a \in \mathcal{A}(s)$}
    \IF{$N(s, a) \ge n_{\mathrm{min}}^{\mathrm{visits}}$}
      \STATE $\mathrm{bonus}(a) \gets c_{\mathrm{puct}}\, P(a \mid s)\,
                              \sqrt{N(s)} / (1 + N(s, a))$
      \STATE $\mathrm{bonus}_{\mathrm{var}}(a) \gets
                  -c_{\mathrm{var}}\, \sigma(s, a) / (1 + N(s, a))$
      \STATE $\mathrm{score}(a) \gets Q(s, a) + \mathrm{bonus}(a) +
                                    \mathrm{bonus}_{\mathrm{var}}(a)$
    \ELSE
      \STATE $\mathrm{score}(a) \gets Q(s, a) + c_{\mathrm{puct}}\,
                                    P(a \mid s)\,
                                    \sqrt{N(s)}$
    \ENDIF
  \ENDFOR
  \STATE $a^{\star} \gets \arg\max_{a} \mathrm{score}(a)$
  \STATE \textbf{return} $a^{\star}$
\end{algorithmic}

\paragraph{Complexity.}
Per-node selection $O(|\mathcal{A}(s)|)$ dominated by score
evaluation; variance penalty adds one division / comparison per
arm and is amortised $O(1)$.  When $\sigma(s, a) = 0$ the
variance term vanishes exactly and the canonical PUCT rule
(Silver 2016 + Rosin 2011) is recovered.

\paragraph{Integration point.}
Wf-11 \texttt{VarianceControlledPUCT} wraps
\texttt{MCTSProofSearch} (R17 ship) with a node-centric selector;
TODO-30 P5.2 \texttt{VarianceMCTSWrapper} exposes
\texttt{search} / \texttt{tree\_statistics} so callers can attach
variance control to the R17 MCTS without patching its internals
(read-mostly on the underlying MCTS).

% ===========================================================================
\paragraph{Summary.}
Six algorithms M22, M23, M24, M25, M26, M28.  M22--M24 are the
three R18 math adapters (W2.1, W2.2, W2.3); M25 is the unifying
OTPD frame; M26, M28 are the Score-Repellent MC and Variance-PUCT
MCTS primitives consumed by the SynFlowNet inner loop.  Each is
CPU-only on the math layer and ships with a W2.x adapter for the
live pipeline.

\endinput